\section{A New Scaling Dimension}

Fusion creates a second way for an AI system to improve. The first is familiar:
better data, algorithms, and compute produce more capable models. We call this
\emph{supply-side intelligence improvement}. It remains indispensable, and
Fusion is designed to inherit every advance it creates.

The second is \emph{demand-side allocation intelligence}: improving how well a
system discovers what a task needs and converts the available model portfolio
into useful evidence and action. The two forms of progress are complementary:

\begin{equation}
  \text{system progress}
  \quad\longleftarrow\quad
  \underbrace{\text{model progress}}_{\text{supply}}
  \quad\times\quad
  \underbrace{\text{allocation progress}}_{\text{demand}}.
\end{equation}

The expression is not a proposed power law. It captures a strategic property
of the system. A stronger model improves the intelligence available to Fusion.
A better task representation, capability map, acquisition policy, or stopping
rule improves the value extracted from every model already available. When
both advance, their effects can compound.

\subsection{Capability without proportional frontier compute}

This scaling dimension is visible in the relationship between performance and
compute. A system that understands its information state should know when
frontier judgment is load-bearing, when a specialist or efficient model can
supply the missing evidence, when only a tool can resolve the uncertainty, and
when the decision no longer benefits from another call.

Compute efficiency is therefore not a purchasing optimization applied after
quality. It is an external sign of allocation intelligence. Redundant frontier
calls indicate that the system cannot distinguish available evidence from
missing evidence. Premature use of a weak source indicates that it cannot
recognize where stronger judgment is necessary. Fusion must improve both sides
of this boundary.

The target is a better performance--compute frontier: at comparable quality,
less frontier computation; at comparable compute, greater task capability.
This is categorically different from reducing cost by accepting weaker
outcomes.

\subsection{Asymmetry is a consequence, not the thesis}

Fusion often selects an asymmetric architecture because the work inside a task
is itself asymmetric. Search, construction, critique, verification, and final
judgment do not necessarily require the same model. But heterogeneous supply
does not create value on its own. A weak or correlated source that contributes
misleading evidence is wasteful at any price
\citep{li2025rethinkingmoa}.

The underlying advantage is the allocation layer: task-state understanding,
conditional capability knowledge, role assignment, and stopping. A system can
select another frontier model when its expected evidence value warrants the
cost. It can select an efficient specialist when that source is more
informative. And it can select no model when the next uncertainty is already
resolved or requires an environmental observation.

\subsection{Scaling through learning}

Allocation intelligence begins with explicit system design, but it can improve
from experience. Each completed task reveals something about which uncertainty
was consequential, which source changed the decision, which evidence was
redundant, and whether the resulting commitment succeeded. That history can
train better task representations, capability estimates, acquisition policies,
and stopping decisions.

Fusion therefore has two independent sources of forward progress. It improves
when models improve, and it improves again when the system learns how to use
them.

\principle{Fusion aims to make system capability compound faster than
frontier-compute consumption.}
